% Part: sets-functions-relations
% Chapter: sets
% Section: proofs-about-sets

\documentclass[../../../include/open-logic-section]{subfiles}

\begin{document}

\olfileid{sfr}{set}{prf}
\olsection{Bukti ngenani Himpunan}

\begin{editorial}
Perangan iki wis diganti dening \verb|content/method/proofs| lan
ora kalebu ing bab iki kanthi baku.
\end{editorial}

\begin{explain}
Himpunan lan notasi kang wis ditepangake menehi cara ringkes kang
migunani kanggo nemtokake himpunan lan ngandharake sesambungane.
Asring kita uga perlu mbuktekake pratelan ngenani sesambungan kasebut.
Yen sampeyan durung kulina karo bukti matematika, iki bisa dadi
prakara anyar. Mula kita bakal njlentrehake conto prasaja.
Kita bakal mbuktekake yen kanggo sembarang himpunan $X$ lan $Y$,
tansah lumaku $X \cap (X \cup Y) = X$. Kepiye mbuktekake
identitas antarane himpunan kaya mangkono? Elinga yen himpunan mung
ditemtokake dening !!{element}s, yaiku himpunan padha yen lan mung
yen nduweni !!{element}s kang padha. Mula kita kudu mbuktekake yen
(a) saben !!{element} saka $X \cap (X \cup Y)$ uga minangka
!!{element} saka $X$ lan, kosok baline, (b) saben !!{element}
saka $X$ uga minangka !!{element} saka $X \cap (X \cup Y)$.
Tegese, kita nuduhake (a) $X \cap (X \cup Y) \subseteq X$ lan
(b) $X \subseteq X \cap (X \cup Y)$.

Pratelan umum kayata ``saben !!{element}~$z$ saka $X \cap (X
\cup Y)$ uga minangka !!{element} saka $X$'' dibuktekake kanthi
nganggep dhisik yen ana sembarang $z \in X \cap (X \cup Y)$,
banjur mbuktekake saka anggepan iki yen $z \in X$.
Sampeyan bisa uga ngerti pola iki minangka ``bukti kondisional umum.''
Ing bukti iki kita uga kudu nggunakake definisi ing anggepan lan
kesimpulan, contone definisi ``$\cap$'' lan ``$\cup$.''
Mula kasus (a) bisa diandharake kaya mangkene:

\begin{quote}
(a) Dhisik kita arep nuduhake $X \cap (X \cup Y) \subseteq X$,
yaiku miturut definisi $\subseteq$, yen $z \in X \cap (X \cup Y)$,
mula $z \in X$, kanggo sembarang~$z$. Dadi anggepen
$z \in X \cap (X \cup Y)$. Amarga $z$ minangka !!{element} saka
irisan loro himpunan yen lan mung yen minangka !!{element} saka
kalorone, kita bisa nyimpulake $z \in X$ lan uga $z \in X \cup Y$.
Mligine, $z \in X$, yaiku kang arep kita tuduhake.
\end{quote}

Iki ngrampungake separo kapisan saka bukti. Ing langkah pungkasan,
kita nggunakake kasunyatan yen sawijining konjungsi ($z \in X$ lan
$z \in X \cup Y$) asale saka anggepan, mula saben konjung uga
asale saka anggepan kang padha. Sampeyan bisa uga ngerti aturan iki
minangka ``eliminasi konjungsi,'' utawa eliminasi $\land$.
Saiki ayo mbuktekake (b):

\begin{quote}
(b) Saiki kita mbuktekake $X \subseteq X \cap (X \cup Y)$,
yaiku miturut definisi $\subseteq$, yen $z \in X$, mula uga
$z \in X \cap (X \cup Y)$, kanggo sembarang~$z$.
Anggepen $z \in X$. Kanggo nuduhake $z \in X \cap (X \cup Y)$,
kita kudu nuduhake (miturut definisi ``$\cap$'') yen (i) $z \in X$
lan uga (ii) $z \in X \cup Y$. Ing kene (i) mung anggepan kita,
dadi ora ana maneh kang kudu dibuktekake. Kanggo (ii), elinga yen $z$
minangka !!{element} saka gabungan himpunan yen lan mung yen
minangka !!{element} saka paling ora salah siji himpunan kasebut.
Amarga $z \in X$, lan $X \cup Y$ minangka gabungan $X$ lan $Y$,
syarat iki lumaku. Dadi $z \in X \cup Y$. Kita wis nuduhake
(i) $z \in X$ lan (ii) $z \in X \cup Y$, mula miturut definisi
``$\cap$,'' $z \in X \cap (X \cup Y)$.
\end{quote}

Andharan iki rada dawa, nanging nuduhake carane kita mikir ngenani
himpunan lan sesambungane. Biasane kita ora ngrinci kaya mangkene;
mligine, kita bisa ora mbaleni kabeh definisi. Bukti asil iki ing
wacan kang luwih lanjut bisa luwih ringkes. Wujude bisa kaya mangkene.
\end{explain}

\begin{prop}[Panyerepan]
Kanggo kabeh himpunan $X$, $Y$,
\[
X \cap (X \cup Y) = X
\]
\end{prop}

\begin{proof}
(a) Upamane $z \in X \cap (X \cup Y)$. Mula $z \in X$, dadi
$X \cap (X \cup Y) \subseteq X$.

(b) Saiki upamane $z \in X$. Mula uga $z \in X \cup Y$,
lan saka iku uga $z \in X \cap (X \cup Y)$.
Dadi $X \subseteq X \cap (X \cup Y)$.
\end{proof}

\begin{prob}
Buktekna kanthi rinci yen $X \cup (X \cap Y) = X$.
Banjur wenehana bukti kang luwih ringkes. (Pituduh: kanggo arah
$X \cup (X \cap Y) \subseteq X$, sampeyan perlu bukti kanthi
kasus, uga diarani eliminasi $\lor$.)
\end{prob}

\end{document}
